% 将本段按论文版式调整后放入附录。路径与提交候选包保持一致。
\section{支撑材料文件列表}
本论文的支撑材料按模型源码、本地实验、官方演练和正式测试四类组织：
\begin{enumerate}
  \item \texttt{程序源码/问题1\_2\_几何与选点/}：问题一、二的几何与定位选点源码、测试及模型说明；
  \item \texttt{程序源码/问题3\_论文方法\_optical/}：问题三论文方法及本地重建模拟入口；
  \item \texttt{程序源码/问题3\_正式测试\_v3\_origin20/}：问题三正式测试使用的完整入口与源码；
  \item \texttt{程序源码/问题4\_V6\_Lite/}：问题四 V6 Lite 的完整入口、模型与源码；
  \item \texttt{本地实验/}：问题三 2000 个共享场景的五方法配对结果，以及问题四 1140 个场景、3420 次运行的配对验证；
  \item \texttt{官方演练/问题3\_问题4\_各100次/}：两题各 100 次官方演练的逐案例总表和独立审计汇总；
  \item \texttt{正式测试/}：问题三、四各三轮的原始加密日志、匿名化审计摘要与逐步过程；
  \item \texttt{MANIFEST.sha256} 与 \texttt{文件清单.tsv}：全包文件校验及目录索引。
\end{enumerate}
其中，本地重建模拟仅用于方法比较；官方演练与正式测试均单独标注版本及证据边界。
